% Saryu, paper B draft v1: what a trained reflection transport does when it cannot solve a group.
% Build:  cd paper && pdflatex paperB && pdflatex paperB
\documentclass[10pt]{article}
\usepackage[margin=1in]{geometry}
\usepackage[T1]{fontenc}
\usepackage{lmodern}
\usepackage{amsmath,amssymb}
\usepackage{booktabs,xcolor,microtype,caption,enumitem,graphicx,float}
\usepackage{tikz,pgfplots}
\pgfplotsset{compat=1.17}
\usepackage[colorlinks=true,linkcolor=blue!50!black,citecolor=blue!50!black,urlcolor=blue!50!black]{hyperref}

\newtheorem{proposition}{Proposition}
\newenvironment{proof}{\par\noindent\textit{Proof.}\ }{\hfill$\square$\par\medskip}
\newcommand{\R}{\mathbb{R}}
\newcommand{\bu}{\mathbf{u}}
\newcommand{\bh}{\mathbf{h}}
\newcommand{\bv}{\mathbf{v}}

\title{\textbf{Quantised failure: trained reflection transports collapse onto exact group quotients}\\[4pt]
\large Technical report --- draft v1}
\author{Varun Sharma}
\date{16 September 2026}

\begin{document}
\maketitle

\begin{abstract}
A recurrent layer whose state is moved by products of Householder reflections can be trained on the
word problem of a finite group: read a sequence of group elements, report the product. When such a
layer succeeds it is exact at any length, because it has learned a homomorphism. When it fails, it
does not fail arbitrarily. Its accuracy lands on $1/|N|$ for a normal subgroup $N$ --- it has learned
the quotient $G/N$, gets the coset right, and guesses inside it. We give the proven ceiling that
makes the phenomenon measurable, the rung structure across several groups, a control at matched group
order that separates the lattice from the order, a table-free diagnostic based on characters, and the
tasks where a quotient is sufficient rather than a failure. We also report the same invariant
appearing in an unrelated role: the normalised trace that identifies which quotient a transport
learned is the quantity that governs how much information the same transport can bind and store.
\end{abstract}

% ======================================================================================
\section{The phenomenon}

Train the transport of Section~\ref{sec:setup} on the word problem of $S_4$ and it either reaches
$1.000$ and stays there as the evaluation length grows to $8\times$ the training length, or it settles
on one of a small set of values: $0.250$, $0.083$, $0.042$. Those are $1/4$, $1/12$, $1/24$, and the
normal subgroups of $S_4$ are $1$, $V_4$ (order 4), $A_4$ (order 12) and $S_4$ itself. A model that
has learned $G/N$ knows which coset the answer lies in and cannot do better than guessing inside it,
so its accuracy is exactly $1/|N|$.

The claim of this paper is that this is what failure looks like for this architecture, that it can be
verified rather than inferred from a coincidence of four numbers, and that the same algebra predicts a
second, apparently unrelated limitation of the same layer.

\paragraph{What makes it checkable.} Three things, in increasing strength.
\begin{enumerate}[leftmargin=*,itemsep=2pt]
  \item \textbf{Flatness.} A homomorphism is length-independent. Every claim below is scored at a
        length far beyond training, and runs whose accuracy drifts with length are excluded rather
        than scored: they are not homomorphisms and the law says nothing about them.
  \item \textbf{Coset consistency.} If the model learned $G/N$, its \emph{errors} must respect
        cosets: predicted and true element in the same coset of $N$, essentially always. A model that
        is merely bad has no reason to respect any coset.
  \item \textbf{A proven ceiling.} For a commuting transport the product is order-independent, so the
        output can only be a function of the multiset of inputs; the best such function is computable
        by enumeration and is a hard bound.
\end{enumerate}

% ======================================================================================
\section{Setup}
\label{sec:setup}

The setup is deliberately not the language model of paper A. There is no input projection and no
gate: each group element $g$ owns a learned table of $n_h$ reflection vectors $\bu_{g,1..n_h}\in\R^d$
and coefficients $\beta_{g,i}$, and the transport it induces is their product,
\[
  T_g \;=\; \prod_{i=1}^{n_h} \bigl(I - \beta_{g,i}\,\hat\bu_{g,i}\hat\bu_{g,i}^\top\bigr),
  \qquad \hat\bu = \bu/\|\bu\| ,
\]
initialised at $\beta=2$, i.e.\ at pure reflections. The state starts at a learned $\bh_0$, each token
applies its own $T_g$, and a linear readout maps the state to the $|G|$ classes. Giving every element
its own transport removes the question of whether a shared projection can separate the elements and
leaves only the question this paper asks: can a \emph{product of reflections} realise the group?

Training reads a word $g_1\dots g_L$ and is supervised on the running product at every prefix, with
an explicit group-law penalty added to the cross-entropy:
\[
  \mathcal{L} \;=\; \mathrm{CE}\bigl(\text{readout}(\bh_t),\, g_1\cdots g_t\bigr)
  \;+\; \lambda\,\mathbb{E}_{\bh,a,b}\,\bigl\|\,T_b T_a \bh - T_{ab}\,\bh\,\bigr\| / \sqrt{d},
  \qquad \lambda = 1 ,
\]
the expectation taken over 64 random triples per step. The penalty asks directly for a homomorphism
rather than hoping one emerges, which is what makes a failure interpretable: a model that ends with a
small penalty and sub-unit accuracy has learned \emph{some} homomorphism, and the question is which.

Unless stated otherwise: $n_h=4$, $d=4$ for $S_4$ and $Q_8$ and $d=5$ for $A_5$, AdamW at
$2\times10^{-3}$ with weight decay $0.01$ and gradients clipped at $1.0$, batch 256, 2000 steps,
trained at length $L=12$ and evaluated at $L=12$ and $L=96$ on 1024 held-out words each. A run counts
as \emph{flat}, and so as a candidate homomorphism, when the two accuracies differ by less than
$0.05$.

% ======================================================================================
\section{The abelian ceiling}
\label{sec:ceiling}

Before asking which quotient a transport learns, it is worth fixing a bound that no amount of
training can move, because it converts one architectural property --- whether the transports commute
--- into a number.

\begin{proposition}[multiset bound]
If the transports commute, the accuracy on the length-$L$ word problem of $G$ cannot exceed
\[
  C(G,L) \;=\; \sum_{M} \Pr[M]\;\max_{g\in G}\,\Pr\bigl[\,g_1\cdots g_L = g \mid M\,\bigr],
\]
the sum running over multisets $M$ of $L$ elements.
\end{proposition}
\begin{proof}
If $T_aT_b=T_bT_a$ for all $a,b$, the state after reading a word depends on the word only through how
many times each element occurs, so the prediction is a function of $M$. The best such function
returns, for each $M$, the most probable product, which is the stated sum.
\end{proof}

The bound is computable by enumeration, and at $L=8$ we evaluate it exactly over all $1287$ multisets
of six elements: $C(\mathbb{Z}_6,8)=1.0000$ and $C(S_3,8)=0.3850$. The pair is matched in order, so a
transport that clears $0.385$ on $S_3$ has done something order-sensitive, and nothing about the size
of the group can explain it away.

\begin{table}[H]\centering\small
\caption{Matched-order test at $L=8$, three seeds. The ceiling is proven, not fitted. Commuting
transports sit on it; the reflection transport passes it on $S_3$.}
\begin{tabular}{llcc}
\toprule
transport & commuting? & $\mathbb{Z}_6$ & $S_3$ \\
\midrule
diagonal rotation & yes & $1.000\pm0.000$ & $0.388\pm0.006$ \\
single Givens     & yes & $1.000\pm0.000$ & $0.389\pm0.003$ \\
Householder product ($n_h=4$) & no & $1.000$ & $0.778$ \\
butterfly         & no  & $1.000\pm0.000$ & $0.779\pm0.380$ \\
transformer       & --- & $0.166\pm0.004$ & $0.167\pm0.003$ \\
\midrule
\textit{proven ceiling for any commuting transport} & & $1.000$ & $\mathbf{0.385}$ \\
\bottomrule
\end{tabular}
\label{tab:ceiling}
\end{table}

Two honest qualifications. The $\pm0.380$ on the butterfly row is a spread across seeds wide enough
that some seeds solve the task and others fail below the ceiling; the mean exceeding $0.385$ says the
ceiling \emph{can} be passed, not that it reliably is. And clearing the ceiling on $S_3$ does not
imply clearing it everywhere: on the same ladder the reflection transport reaches $0.374$ on $Q_8$
against a ceiling of $0.544$, $0.085$ on $A_4$ against $0.286$, and chance on $S_4$ and $A_5$. Being
non-commuting removes a ceiling; it does not supply an optimiser that finds what lies above it. The
rest of this paper is about where those runs land instead.

% ======================================================================================
\section{The lattice law}
\label{sec:lattice}

A transport that ends training with a small group-law penalty is a homomorphism $\rho: G\to O(d)$
composed with a readout. Its kernel $N=\ker\rho$ is a normal subgroup, elements of the same coset of
$N$ are indistinguishable to it, and its accuracy is therefore $1/|N|$. The set of achievable
accuracies is thus the set $\{1/|N| : N \trianglelefteq G\}$ --- the \emph{rungs} --- which is a
property of the group alone and is invisible to the model. We enumerate the normal subgroups of each
group directly (unions of conjugacy classes closed under the product) rather than quoting them:
$S_4$ gives orders $1,4,12,24$; $Q_8$, being Hamiltonian, gives $1,2,4,4,4,8$; $A_5$, being simple,
gives only $1,60$.

\begin{table}[H]\centering\small
\caption{Four seeds per group. Every run is flat, so every run is a candidate homomorphism, and every
one lands on a rung of its own group. Worst deviation across all twelve runs: $0.021$.}
\begin{tabular}{lcccc}
\toprule
group (rungs) & seed & $L=12$ & $L=96$ & nearest rung (dev.) \\
\midrule
$Q_8$  & 0 & 0.247 & 0.241 & 0.250 (0.009) \\
$1.000\;0.500\;0.250\;0.125$ & 1 & 0.498 & 0.521 & 0.500 (0.021) \\
                             & 2 & 1.000 & 1.000 & 1.000 (0.000) \\
                             & 3 & 0.264 & 0.243 & 0.250 (0.007) \\
\midrule
$S_4$  & 0 & 1.000 & 1.000 & 1.000 (0.000) \\
$1.000\;0.250\;0.083\;0.042$ & 1 & 0.081 & 0.100 & 0.083 (0.017) \\
                             & 2 & 0.264 & 0.251 & 0.250 (0.001) \\
                             & 3 & 0.257 & 0.259 & 0.250 (0.009) \\
\midrule
$A_5$  & 0 & 0.015 & 0.020 & 0.017 (0.003) \\
$1.000\;0.017$               & 1 & 0.009 & 0.021 & 0.017 (0.004) \\
                             & 2 & 0.018 & 0.014 & 0.017 (0.003) \\
                             & 3 & 0.015 & 0.016 & 0.017 (0.001) \\
\bottomrule
\end{tabular}
\label{tab:rungs}
\end{table}

\paragraph{Which rungs are actually informative.} The bottom rung $1/|G|$ coincides with chance, so a
run landing there is consistent with the law but tells us nothing: it is the trivial quotient and the
absence of learning at once. This applies to all four $A_5$ runs, at $0.017$. The registered
falsifier for $A_5$ --- a flat accuracy strictly between $0.017$ and $1.000$, which a simple group
forbids --- was not triggered, but neither did any $A_5$ run solve the task, so the group contributes
a consistency check rather than a positive result. The load is carried by the intermediate rungs that
are neither chance nor perfect: $Q_8$ at $0.500$ and $0.250$, and $S_4$ at $0.250$ and $0.083$.

\paragraph{Coset consistency.} Accuracy alone could be a numerical coincidence, so we test the
mechanism on the errors. If a model learned $G/N$, its predictions must land in the correct coset of
$N$ essentially always, and the kernel is identified as the \emph{smallest} $N$ whose coset score
reaches $1.000$ --- read off the errors, with no reference to the accuracy. The accuracy predicted by
that independently identified $N$ then matches the measured one on all five $S_4$ runs tested:

\begin{table}[H]\centering\small
\caption{The subgroup identified by the errors predicts the accuracy. Coset scores are the fraction
of predictions landing in the correct coset.}
\begin{tabular}{cccccll}
\toprule
raw acc. & $1$ & $V_4$ & $A_4$ & $S_4$ & identified $N$ & $1/|N|$ \\
\midrule
1.000 & 1.000 & 1.000 & 1.000 & 1.000 & $1$ (faithful) & 1.000 \\
1.000 & 1.000 & 1.000 & 1.000 & 1.000 & $1$ (faithful) & 1.000 \\
0.230 & 0.230 & 0.999 & 0.999 & 1.000 & $V_4$ & 0.250 \\
0.081 & 0.081 & 0.329 & 1.000 & 1.000 & $A_4$ & 0.083 \\
0.042 & 0.042 & 0.165 & 0.472 & 1.000 & $S_4$ & 0.042 \\
\bottomrule
\end{tabular}
\label{tab:coset}
\end{table}

\paragraph{A measurement we withdraw.} Our earlier runs also reported a per-run kernel size, obtained
by clustering the transports $M[g]=\mathbb{E}_\bh\,T_g\bh$ and setting $|\ker| = |G|/(\#\text{clusters})$.
That estimator is invalid: $T_g$ is linear and $\bh$ is zero-mean, so $M[g]$ is a Monte-Carlo zero for
every $g$ and the clustering runs on sampling noise. Measured, $\|M[g]\|$ is $5.7\%$ of the typical
$\|T_g\bh\|$, below the clustering threshold; applied to an \emph{untrained} model whose $24$
transports are independent random draws and whose true kernel is trivial, it reports $|\ker|=8$. It
returned orders that are not normal subgroup orders at all --- $8$ for $S_4$, $30$ for $A_5$ --- in 4
of 12 runs, which is how the fault was found. Nothing in this paper rests on that column; the kernel
is identified by coset consistency instead.

% ======================================================================================
\section{The same-order control}

The cross-group evidence is confounded by group order, since a 24-element problem and an 8-element
problem differ in chance level as well as in lattice. The control removes it: $S_4$ and $\mathrm{SL}(2,3)$
both have order 24, and their rungs are disjoint apart from the endpoints.
\[
  S_4:\; 1.000,\ 0.250,\ 0.083,\ 0.042 \qquad\qquad \mathrm{SL}(2,3):\; 1.000,\ 0.500,\ 0.125,\ 0.042
\]
Across four registered runs (2400 steps, 1200 steps, an $S_4$-only run at 1800 steps, and a run at
width 3), 28 failures were flat enough to score. Every one landed on a rung of its own group, 13 of
them on a rung the other group does not possess, and none within 0.03 of the other group's exclusive
rungs. One run missed its registered tolerance --- $\mathrm{SL}(2,3)$ at 0.469 against a predicted
0.500, with coset consistency 1.000 --- so the registered verdict is \emph{falsified}, and we report it
as such; in binomial standard errors that miss is $z=-2.00$, the only $|z|>1.96$ among the 28, which
says the tolerance was set too tight rather than that the law failed. The runs, their predictions and
their verdicts are in the repository.

% ======================================================================================
\section{Characters as a table-free diagnostic}
\label{sec:char}

Identifying the kernel by coset consistency needs the group table. There is a diagnostic that does
not. For a representation $\rho$ with character $\chi(g)=\operatorname{tr}\rho(g)$, the norm
\[
  \langle\chi,\chi\rangle \;=\; \frac{1}{|G|}\sum_{g\in G}\chi(g)^2 \;=\; \sum_i m_i^2
\]
counts irrep multiplicities, and needs only traces of the learned transports --- no labels, no Cayley
table, no character table. Enumerating every way to build a $4$-dimensional representation of $S_4$
and intersecting the kernels of its parts gives what each value implies:

\begin{center}\small
\begin{tabular}{lll}
\toprule
kernel & quotient realised & attainable $\langle\chi,\chi\rangle$ at $d=4$ \\
\midrule
$1$ & $S_4$ (faithful) & $2$ \\
$V_4$ & $S_3$ & $3,\;4,\;5$ \\
$A_4$ & $\mathbb{Z}_2$ & $8,\;10,\;16$ \\
$S_4$ & trivial & $16$ \\
\bottomrule
\end{tabular}
\end{center}

At $d=4$ the value $2$ is attainable \emph{only} by a faithful representation, so it is a sufficient
certificate of faithfulness. That statement is conditional on full rank, and the condition is not
decorative: a trained transport may let a dimension die, and at $d=3$ the value $2$ means kernel
$V_4$, at $d=2$ kernel $A_4$. We observe exactly this --- runs that solve the task with
$\langle\chi,\chi\rangle\approx1.12$ and a single $3$-dimensional irrep carrying the whole state. The
diagnostic must therefore be read together with the effective rank, not on its own.

A second condition matters more. $\langle\chi,\chi\rangle$ is a character norm only if the transport
\emph{is} a representation. In our runs the solved models carry a group-law violation of $\approx
0.008$ and measure $\langle\chi,\chi\rangle = 1.99$--$2.01$, agreeing with the table; the failed runs
carry violations of $0.10$--$0.29$, are not homomorphisms at all, and their measured values
($4.4$--$5.5$) are not character norms of anything and should not be read as quotient labels. The
predicted quotient values are so far confirmed in one clean instance rather than as a set: a model
that solves the $V_4$ task exactly (Section~\ref{sec:quotient}) decomposes as two copies of the
trivial irrep plus the $2$-dimensional one, giving $\langle\chi,\chi\rangle = 2^2+1^2 = 5$ with
kernel $V_4$, as the table requires.

Finally, the diagnostic does not work as a training signal. An arm that adds a penalty driving the
mean squared trace toward $2.0$ --- no table, no labels, nothing but traces --- solved 2 of 6 seeds
against 3 of 6 for the unregularised baseline. It certifies a representation; it does not produce
one.

% ======================================================================================
\section{The same invariant governs binding}
\label{sec:binding}

The diagnostic above is a character: $\chi(g)=\operatorname{tr}(g)$, and $\langle\chi,\chi\rangle$
identifies the learned quotient without the group table. The same quantity controls a capability that
has nothing to do with group word problems.

For an orthogonal $g$ acting on the sphere, the expected overlap of a random unit vector with its
image is the normalised trace,
\[
  \mathbb{E}\,[\,\bv^\top g\,\bv\,] = \operatorname{tr}(g)/d ,
\]
which is exactly how much an operation \emph{hides} a vector --- the property a binding operation
needs if a value is to be stored against a key and recovered later. A single Householder reflection
has $\operatorname{tr}=d-2$, the least-hiding non-trivial element of $O(d)$, and a product of $k$
reflections gives overlap $\approx e^{-2k/d}$. Measured at head dimension 64, with weights set by hand
and no training, capacity follows the trace exactly: four associations are recalled at $0.453$ with
$k=2$ and at $1.000$ with $k=32\approx d/2$.

So the same transport is a tracker at small $k$ and an associative store at $k\approx d/2$, and the
character is what separates the regimes. A layer that composes group elements perfectly and cannot
retain one fact among four is not exhibiting two unrelated limitations; it is at one end of a single
curve.

% ======================================================================================
\section{What a quotient is good for}
\label{sec:quotient}

Calling a quotient a failure assumes the task needed the whole group. Change the target to the coset
of $N$ and the same collapse becomes the correct answer: the $V_4$ task asks for the $V_4$-coset of
the product ($6$ classes, chance $0.167$) and needs only $S_4/V_4\cong S_3$; the $A_4$ task asks for
the $A_4$-coset ($2$ classes, chance $0.500$) and needs only $S_4/A_4\cong\mathbb{Z}_2$.

Solve rates rise as the required group shrinks --- 2 of 4 seeds on the full task, 2 of 4 on $V_4$,
3 of 4 on $A_4$ --- but the sharper result is where the \emph{failures} land. If the law of
Section~\ref{sec:lattice} is really about the group being learned, then on a quotient task it should
apply to the quotient, and the rungs should be those of $S_3$ and $\mathbb{Z}_2$ rather than those of
$S_4$. The normal subgroups of $S_3$ have orders $1,3,6$, giving rungs $1.000$, $0.333$, $0.167$;
those of $\mathbb{Z}_2$ give $1.000$ and $0.500$. Neither $0.333$ nor $0.500$ is a rung of $S_4$.

\begin{table}[H]\centering\small
\caption{The lattice law one level down. All eight runs land on a rung of the \emph{quotient} they
are asked for; worst deviation $0.013$.}
\begin{tabular}{llcccc}
\toprule
task & group needed & \multicolumn{4}{c}{$L=96$ accuracy by seed} \\
\midrule
$V_4$-coset & $S_3$ (rungs $1.000\;0.333\;0.167$) & 1.000 & 1.000 & 0.320 & 0.336 \\
$A_4$-coset & $\mathbb{Z}_2$ (rungs $1.000\;0.500$) & 1.000 & 0.495 & 1.000 & 1.000 \\
\bottomrule
\end{tabular}
\label{tab:quotient}
\end{table}

The failures sit at $0.320$ and $0.336$ against $S_3$'s rung of $1/3$, and at $0.495$ against
$\mathbb{Z}_2$'s rung of $1/2$. A registered prediction accompanies this: a quotient task does not
require a faithful representation, so the amplitude on $S_4$'s $3$-dimensional irreps should stay near
zero even when the task is solved perfectly. It does --- at most $0.042$ across all eight quotient
runs, against $\approx1.00$ on the solved full-task runs. Capacity tracks the task, not the ambient
group.

% ======================================================================================
\section{Limitations}

Small groups, small dimensions, single-layer transports, four seeds per cell, and a training
objective that includes an explicit group-law penalty --- the law describes where this architecture
lands when it fails, not a claim about trained networks in general. The same-order control is
registered as falsified on one run of 28. The Frobenius--Schur account of \emph{which} rung is
reached, present in earlier notes, is not claimed here: $\mathrm{SL}(2,3)$ solved the task outright at
width 6, which that account does not predict.

Three weaknesses are specific. The bottom rung is chance, so the four $A_5$ runs confirm the law
without demonstrating it, and the positive evidence rests on the intermediate rungs of $Q_8$ and
$S_4$ and on coset consistency. Clearing the abelian ceiling is shown on $S_3$ only; on larger groups
the same transport falls below what a commuting transport achieves, so Section~\ref{sec:ceiling}
bounds what the architecture can express and says nothing about what training reliably finds. And the
character diagnostic is doubly conditional --- on the transport being a homomorphism, and on the
representation being full rank --- with the predicted quotient values so far confirmed in a single
clean instance rather than across the table.

Every derived quantity above --- the ceiling, the rung sets, the attainable character norms --- is
recomputed from the group definitions by \texttt{evidence/verify\_paperB.py}, which also re-reads the
measured tables from the result files rather than trusting a transcription.

\end{document}
